\chapter{Sprint 2 : Module de Découverte et de Mise en Relation}
\minitoc
\newpage

\section{Sprint II : Module de Découverte et de Mise en Relation}

\subsection{Objectif du Sprint}
La valeur centrale de GOSPORT réside dans la capacité à connecter des athlètes avec des professionnels du sport (coachs, nutritionnistes) de manière fluide et sécurisée. Ce deuxième sprint implémente le cycle complet de mise en relation : la découverte de profils, la consultation des fiches détaillées, et le flux de demandes de coaching incluant leur acceptation ou leur refus.

À l'issue de ce sprint, un athlète peut parcourir l'annuaire des professionnels disponibles sur la plateforme, consulter leurs profils, et envoyer une demande de suivi. Le professionnel peut alors consulter les demandes en attente et décider de les accepter ou de les refuser.

\subsection{Backlog du Sprint II}

\begin{table}[H]
\centering
\renewcommand{\arraystretch}{1.4}
\resizebox{\textwidth}{!}{%
\begin{tabular}{|c|c|p{5cm}|p{6.5cm}|c|}
\hline
\rowcolor[HTML]{DAE8FC}
\textbf{ID} & \textbf{Acteur} & \textbf{User Story} & \textbf{Tâche Technique} & \textbf{Priorité} \\ \hline
2.1 & Athlète & Découvrir et filtrer les coachs disponibles & Filtres par spécialité, tarif et note & Must \\ \hline
2.2 & Tous & Consulter le profil public d'un professionnel & Page de détail du profil professionnel & Must \\ \hline
2.3 & Athlète & Envoyer une demande de coaching & Entité CoachingRequest et workflow de demande & Must \\ \hline
2.4 & Coach & Accepter ou refuser une demande & Mise à jour du statut de la demande & Must \\ \hline
2.5 & Athlète & Envoyer une demande à un nutritionniste & Workflow de demande nutritionnelle & Should \\ \hline
2.6 & Nutritionniste & Répondre à une demande de suivi nutrition & Mise à jour du statut nutritionnel & Should \\ \hline
2.7 & Coach & Dissocier un athlète de son suivi & Action de dissociation sécurisée & Could \\ \hline
\end{tabular}%
}
\caption{Backlog du Sprint 2 -- Module de mise en relation}
\end{table}

\section{Analyse}

\subsection{Diagramme de cas d'utilisation du deuxième sprint}
\begin{figure}[H]
\centering
\includegraphics[width=0.88\textwidth]{Sprint2_Cas_Utilisation.png}
\caption{Diagramme de cas d'utilisation -- Sprint 2}
\end{figure}

\subsection{Tableau descriptif du cas d'utilisation « ... »}

\subsection{Tableau descriptif du cas d'utilisation « ... »}

\subsection{Tableau descriptif du cas d'utilisation « ... »}

\subsection{Tableau descriptif du cas d'utilisation « ... »}

\subsection{Tableau descriptif du cas d'utilisation « ... »}

\subsection{Tableau descriptif du cas d'utilisation « ... »}

\subsection{Tableau descriptif du cas d'utilisation « ... »}

\subsection{Tableau descriptif du cas d'utilisation « ... »}

\section{Conception}

\subsection{Diagrammes de séquences}

\subsection{A) Diagramme de séquence « ... »}
\begin{figure}[H]
\centering
\includegraphics[width=0.88\textwidth]{Sprint2_Sequence_CoachingRequest.png}
\caption{Diagramme de séquence -- Cycle de vie d'une demande de coaching}
\end{figure}

\subsection{B) Diagramme de séquence « ... »}

\subsection{C) Diagramme de séquence « ... »}

\subsection{D) Diagramme de séquence « ... »}

\subsection{Diagramme de classes du deuxième sprint}
\begin{figure}[H]
\centering
\includegraphics[width=0.9\textwidth]{Sprint2_Classes.png}
\caption{Diagramme de classes -- Mise en relation}
\end{figure}

Les entités clés de ce sprint sont :\texttt{CoachingRequest} (statut: PENDING \textbar{} ACCEPTED \textbar{} REJECTED) et \texttt{NutritionConnection}. Ces deux entités établissent un lien entre un \texttt{User} (athlète) et un professionnel (coach ou nutritionniste). Le cycle de vie de cet objet -- de sa création à son acceptation -- est central pour ouvrir l'accès aux fonctionnalités suivantes (programmes, nutrition, messagerie).

\section{Réalisation}

\subsection{Description des interfaces}

\section*{Conclusion du Sprint II}
\addcontentsline{toc}{section}{Conclusion du Sprint II}
\mtcaddsection
\markboth{Conclusion du Sprint II}{}
Le module de mise en relation est complet. Les athlètes peuvent désormais explorer la place de marché des professionnels, consulter les profils publics et initier des demandes de coaching ou nutritionnelles. Les coachs disposent d'un tableau de bord pour gérer ces demandes entrantes. L'établissement de ces liens est la clé qui déverrouille l'accès aux fonctionnalités des sprints suivants (partage de programmes, nutrition, messagerie).
